\section{Methodology}
\label{methodology}

\subsection{Framework Overview}
\label{method:overview}

The proposed framework verifies the feasibility of natural language requirements before downstream development activities.
Given a requirement statement and available project knowledge, the workflow produces a feasibility judgment, supporting constraint evidence, detected realism issues, and suggested requirement refinements.

\paragraph{Scope.}
This study focuses only on G6-feasibility: requirement feasibility analysis and realism verification.
It does not address requirement graph construction, conflict detection, architecture mapping, executable requirements, requirement planning, or dynamic requirement clarification.

\subsection{Multi-Agent Feasibility Analysis Workflow}
\label{method:workflow}

\subsubsection{Analyst Agent}
\label{method:analyst_agent}

The analyst agent normalizes the input requirement, identifies the intended capability, and extracts feasibility-relevant assumptions.
Its output is a structured requirement interpretation used by the downstream verification agents.

\subsubsection{Constraint Verification Agent}
\label{method:constraint_agent}

The constraint verification agent checks whether the requirement contains explicit or implicit constraints that may exceed project limits.
It focuses on constraint categories rather than broad requirement quality issues.

\subsubsection{Deployment and Resource Agent}
\label{method:deployment_agent}

The deployment and resource agent inspects operational assumptions such as infrastructure availability, runtime environment, data access, hardware limits, and staffing or budget constraints.
It reports whether the requirement appears deployable under the available context.

\subsection{Requirement Feasibility Detection}
\label{method:detection}

\subsubsection{Resource Constraints}
\label{method:resource_constraints}

Resource constraints include budget, staffing, hardware capacity, data availability, and external service dependencies.

\subsubsection{Performance Constraints}
\label{method:performance_constraints}

Performance constraints include latency, throughput, scalability, reliability, and accuracy targets that may be unrealistic for the stated system context.

\subsubsection{Deployment Constraints}
\label{method:deployment_constraints}

Deployment constraints include platform compatibility, network availability, privacy restrictions, operating environments, and integration conditions.

\subsubsection{Temporal Constraints}
\label{method:temporal_constraints}

Temporal constraints include deadlines, real-time behavior, release schedules, and maintenance windows that may conflict with implementation feasibility.

\subsection{Knowledge-Guided Constraint Analysis}
\label{method:knowledge}

\subsubsection{Requirement Knowledge}
\label{method:requirement_knowledge}

Requirement knowledge includes domain terms, stakeholder assumptions, requirement templates, and known feasibility risks.

\subsubsection{Constraint Rules}
\label{method:constraint_rules}

Constraint rules encode feasibility checks such as resource budget limits, deployment prerequisites, and incompatible constraint combinations.

\subsubsection{Domain Knowledge}
\label{method:domain_knowledge}

Domain knowledge provides contextual evidence for judging whether a requirement is realistic in a target application domain.

\subsubsection{Prompt Strategies}
\label{method:prompt_strategies}

Prompt strategies guide agents to provide structured judgments, constraint evidence, uncertainty labels, and refinement suggestions.

\subsection{Requirement Verification and Refinement}
\label{method:verification_refinement}

\subsubsection{Feasibility Verification}
\label{method:feasibility_verification}

The final feasibility decision aggregates agent findings into labels such as feasible, conditionally feasible, infeasible, or insufficient information.

\subsubsection{Unrealistic Requirement Detection}
\label{method:unrealistic_detection}

Unrealistic requirements are flagged when stated goals conflict with known constraints, require unavailable resources, or assume unsupported deployment conditions.

\subsubsection{Requirement Adjustment Suggestions}
\label{method:adjustment}

For infeasible or conditionally feasible requirements, the workflow suggests conservative refinements such as relaxing targets, adding missing assumptions, or replacing unverifiable claims with measurable constraints.

\subsection{Human-in-the-Loop Interaction}
\label{method:hitl}

Human reviewers inspect the agents' feasibility judgments, confirm domain assumptions, and decide whether suggested refinements should be accepted.
The human-in-the-loop step is designed to reduce review effort while preserving stakeholder authority over requirement decisions.
